\documentclass[12pt,a4paper]{article}

\usepackage[utf8]{inputenc}
\usepackage[T1]{fontenc}
\usepackage[brazil]{babel}
\usepackage{geometry}
\usepackage{enumitem}
\usepackage{setspace}

\geometry{margin=2.5cm}

\title{\textbf{Curso Básico de R para Economistas}\\
	Fundamentos}
\author{}
\date{}

\begin{document}
	
	\maketitle
	
	\section*{Curso de Extensão}
	
	\noindent \textbf{Carga horária:} 20h (5 encontros)
	
	\noindent \textbf{Coordenador:} Sérgio Luiz de Medeiros Rivero
	
	\noindent \textbf{Instrutores:} Nathan Pereira Gomes de Sousa, Ryan Pereira Oliveira, Witiny Leandro Pereira Ramos
	
	\noindent \textbf{Público-alvo:} Estudantes de graduação em Economia e áreas afins.
	
	\noindent \textbf{Pré-requisitos:} Não há pré-requisitos formais. Recomenda-se familiaridade básica com regressão linear e conceitos introdutórios de estatística e econometria.
	
	\noindent \textbf{Formato:} 5 encontros presenciais totalizando 20h, combinando exposição teórica e atividades práticas no RStudio.
	
	
	\section*{Objetivo Geral}
	
	
	Introduzir o uso da linguagem R como ferramenta aplicada à Economia, apresentando o fluxo básico de uma análise empírica, desde a organização dos dados até a estimação e interpretação de modelos econométricos simples.
	
	\section*{Objetivos Específicos}
	
	\begin{itemize}
		\item Apresentar o R como ambiente computacional utilizado na pesquisa econômica;
		
		\item Familiarizar os participantes com o RStudio e com boas práticas iniciais de organização e reprodutibilidade;
		
		\item Introduzir noções básicas de programação em R voltadas à análise de dados;
		
		\item Ensinar procedimentos de importação, organização e exploração de bases de dados;
		
		\item Desenvolver visualizações básicas para análise de dados;
		
		\item Estimar e interpretar modelos de regressão linear simples no R.
	\end{itemize}
	
	\section*{Metodologia}
	
	
	O curso será conduzido por meio de aulas expositivas integradas a demonstrações práticas no RStudio. Os participantes acompanharão a execução dos códigos e realizarão exercícios guiados, com ênfase na conexão entre teoria econômica, econometria e implementação computacional.
		
	Será adotada uma abordagem progressiva, iniciando pela organização do ambiente computacional e avançando para aplicações práticas de análise empírica utilizando bases de dados econômicos.
	
	Ao longo das aulas serão introduzidos novos métodos, pacotes e ferramentas conforme o avanço da turma, permitindo contato inicial com recursos modernos do ecossistema R.
	
	\section*{Avaliação}
	
	A avaliação ocorrerá de forma prática e processual, considerando:
	
	\begin{itemize}
		\item Participação nas atividades desenvolvidas durante os encontros;
		
		\item Execução dos exercícios propostos;
		
		\item Elaboração e apresentação final de uma pesquisa empírica simples em script R, contendo procedimentos de importação, tratamento, análise descritiva e/ou estimação de um modelo de regressão linear simples.
	\end{itemize}
	
	As apresentações finais ocorrerão no dia \textbf{19} de junho.
	
	Ao final do curso, os participantes terão acesso aos scripts utilizados durante as aulas, contendo comentários explicativos para estudo posterior.
	
	\section*{Conteúdo Programático}
	
	\subsection*{Encontro 1 - Introdução ao R e ao ambiente computacional}
	
	\begin{enumerate}[label=\arabic*.]
		
		\item Apresentação do curso
		\begin{itemize}
			\item Objetivos e estrutura geral do curso
			\item Fluxo básico de uma análise empírica
			\item O papel do R na pesquisa econômica
		\end{itemize}
		
		\item Introdução ao ambiente computacional
		\begin{itemize}
			\item O que é o R
			\item Diferença entre R e RStudio
			\item Instalação do R e RStudio
			\item Estrutura do ambiente RStudio
		\end{itemize}
		
		\item Organização e reprodutibilidade
		\begin{itemize}
			\item Scripts e projetos no RStudio
			\item Diretórios de trabalho
			\item Organização de arquivos
			\item Boas práticas e erros frequentes
		\end{itemize}
		
	\end{enumerate}
	
	\subsection*{Encontro 2 - Fundamentos do R e importação de dados}
	
	\begin{enumerate}[label=\arabic*.]
		\setcounter{enumi}{3}
		
		\item Conceitos básicos da linguagem R
		\begin{itemize}
			\item Objetos e funções
			\item Vetores e data frames
			\item Noções básicas de sintaxe
			\item Operações básicas no R
		\end{itemize}
		
		\item Importação e exploração de dados
		\begin{itemize}
			\item Conceito de base de dados em Econometria
			\item Importação de arquivos CSV e Excel
			\item Estrutura e visualização de bases de dados
			\item Primeira exploração dos dados
		\end{itemize}
		
	\end{enumerate}
	
	\subsection*{Encontro 3 - Transformação e análise exploratória de dados}
	
	\begin{enumerate}[label=\arabic*.]
		\setcounter{enumi}{5}
		
		\item Transformação de dados
		\begin{itemize}
			\item Seleção e filtragem de dados
			\item Criação de variáveis
			\item Organização e agregação de informações
			\item Introdução ao pacote \texttt{dplyr}
		\end{itemize}
		
		\item Análise exploratória de dados
		\begin{itemize}
			\item Estatísticas descritivas
			\item Interpretação inicial dos dados
			\item Construção de tabelas
			\item Introdução ao pacote \texttt{flextable}
		\end{itemize}
		
	\end{enumerate}
	
	\subsection*{Encontro 4 - Visualização de dados, introdução à econometria no R e comunicação empírica}
	
	\begin{enumerate}[label=\arabic*.]
		\setcounter{enumi}{7}
		
		\item Visualização de dados no R
		\begin{itemize}
			\item Introdução ao pacote \texttt{ggplot2}
			\item Gráficos de barras, linhas e dispersão
			\item Visualização aplicada à análise econômica
			\item Comunicação de resultados empíricos
		\end{itemize}
		
		\item Introdução à regressão linear no R
		\begin{itemize}
			\item Estimação de modelos de regressão linear simples
			\item Interpretação dos coeficientes
			\item Relação entre teoria econométrica e implementação computacional
			\item Apresentação de resultados
		\end{itemize}
		
		\item Apoio ao desenvolvimento da atividade final
		\begin{itemize}
			\item Estruturação do script final
			\item Organização da análise empírica
			\item Revisão dos principais conceitos do curso
		\end{itemize}
		
	\end{enumerate}
	
	\subsection*{Encontro 5 - Avaliação final}
	
	\begin{enumerate}[label=\arabic*.]
		\setcounter{enumi}{10}
		
		\item Avaliação final
		\begin{itemize}
			\item Apresentação da atividade empírica desenvolvida
			\item Discussão dos resultados obtidos
			\item Encerramento e orientações para continuidade dos estudos
		\end{itemize}
		
	\end{enumerate}
	
	\section*{Resultados de Aprendizagem}
	
	Ao final do curso, espera-se que os participantes sejam capazes de:
	
	\begin{itemize}
		\item Utilizar o RStudio como ambiente básico de análise empírica;
		
		\item Importar e organizar bases de dados simples no R;
		
		\item Realizar análises descritivas iniciais;
		
		\item Construir visualizações básicas;
		
		\item Estimar e interpretar modelos de regressão linear simples;
		
		\item Desenvolver scripts reproduzíveis aplicados a problemas econômicos.
	\end{itemize}
	
	\section*{Software e Dados}
	
	\textbf{Software:}
	\begin{itemize}
		\item R
		\item RStudio
	\end{itemize}
	
	\textbf{Pacotes:}
	\begin{itemize}
		\item tidyverse
		\item ggplot2
		\item flextable
	\end{itemize}
	
	\textbf{Fontes de dados:}
	\begin{itemize}
		\item Bases fornecidas durante o curso;
		\item IBGE;
		\item IPEADATA;
		\item Banco Central do Brasil.
	\end{itemize}
	
	\section*{Bibliografia Básica}
	
	\begin{itemize}
		\item WICKHAM, Hadley; ÇETINKAYA-RUNDEL, Mine; GROLEMUND, Garrett. \textit{R para Ciência de Dados}. 2. ed.
		
		\item WOOLDRIDGE, Jeffrey M. \textit{Introdução à Econometria: uma abordagem moderna}.
		
		\item GUJARATI, Damodar N.; PORTER, Dawn C. \textit{Econometria Básica}.
	\end{itemize}
	
\end{document}